The cylinder question belongs to the line of covering problems that begins
with Tarski's plank problem. A \emph{plank} is the region between two
parallel hyperplanes, and its width is their distance. Bang proved that
a finite plank covering of a convex body has total width at least the
body's minimum width~\cite{Bang1951}; see also
\cite[Theorem~2.1]{Bezdek2009}.
Ball proved the directionwise refinement for centrally symmetric bodies:
after dividing each plank's width by the body's width in the same normal
direction, the sum is at least one~\cite[arXiv v1, Theorem~1]{Ball1991}.
Thus the width cost of a plank cover has both an absolute formulation
and a relative formulation adapted to the individual directions.

For cylinders in three dimensions, the two-cylinder cover of a regular
tetrahedron shows that replacing widths by base areas cannot preserve
the constant one. Its two axes are parallel to opposite edges, and its
total base area equals half the tetrahedron's minimum projection area.
Bezdek~\cite[Problem~3.1]{Bezdek2009} and
Bezdek and Litvak~\cite[arXiv v1, Introduction]{BL2009} attribute this example and
the resulting half-area question to Bang's 1951 paper. Their formulations
allow measurable bases perpendicular to the axes, as in our definition.

For a finite covering $\mathcal C=(B_i+\R u_i)_{i=1}^m$, define the
\emph{directionwise relative area} by
\begin{equation}\label{eq:relative-cost}
 \mathcal R(\mathcal C;K):=
 \sum_{i=1}^m\frac{|B_i|}{|\pi_{u_i^\perp}K|}.
\end{equation}
All denominators are positive because $K$ has nonempty interior.
Bezdek and Litvak proved $\mathcal R(\mathcal C;K)\ge1/3$
for every finite covering in three dimensions, and the stronger bound
$\mathcal R(\mathcal C;K)\ge1$ when $K$ is an
ellipsoid. In their proof, the general case combines the Rogers--Shephard
section--projection inequality with a volume comparison. For ellipsoids,
they reduce to a ball and use a density whose integral along every
interior chord is the same~\cite[arXiv v1, Theorem~3.1 and its proof]{BL2009}.
These arguments leave
a gap between the universal one-third lower bound and the proposed
one-half threshold.

Bezdek and Khan record the directionwise assertion
$\mathcal R(\mathcal C;K)\ge1/2$ under the name
\emph{1-Codimensional Cylinder Covering Conjecture}
\cite[arXiv v2, Definition~3 and Conjecture~4.13]{BK2016}.
It would imply \eqref{eq:conjectured-bound}, since
\begin{equation}\label{eq:cost-comparison}
 \mathcal R(\mathcal C;K)
 \le \frac{\sum_i|B_i|}{A_{\min}(K)}.
\end{equation}
Verreault's 2026 survey records the half-area question as
unresolved~\cite[Question~4.14]{Verreault2026}.
Our construction starts with the tetrahedral equality example itself
and replaces its two axes by finitely many nearby directions. It
violates both half-area bounds while retaining compact triangular bases.
